\section{Discussion}\label{discuss}

\subsection{What the probability diagnostic adds}

The main contribution is not a new HVAC method. It is a diagnostic interface for reading a calibrated TSV probability distribution. Expected TSV, defined as the probability-weighted mean $\mu_{\mathrm{TSV}}=\sum_k k p_k$, remains useful: in the full panel it was strongly associated with $p_{\mathrm{tail}}$, and the diagnostic did not produce near-neutral high-tail states under $|\mu_{\mathrm{TSV}}|<0.15$. This is an important constraint on the interpretation. The probability representation does not overturn the scalar mean.

The added value is that the distribution makes tail probability explicit. A scalar mean reports central tendency; $p_{\mathrm{tail}}$ reports the probability mass in the severe cold and warm categories. Similar mean states can still have different tail probabilities, and the threshold status of $p_{\mathrm{tail}}$ is directly available only when the distribution is retained. In this fixed-reference future-weather panel, the visible tail was also strongly warm-skewed: cold-tail probability was small, and high-tail records were warm-dominant. That pattern is conditional on the selected cities, weather files, fixed thermostat schedule, and static comfort model, but it is exactly the kind of directional diagnostic that the probability object makes explicit.

The panel summary in Table~\ref{tab:panel_tail_summary} also matters because it separates broad stress-test gradients from local case differences. High-tail exceedance increased from the baseline-2020s slice to the late-2080s slice and was higher under SSP5-8.5 than SSP2-4.5, but the hot and heatwave-extreme annual files were close in aggregate exceedance. The city spread was larger than either scenario contrast or annual-severity contrast. This supports a cautious reading: the fixed archetype shows more high-tail states under later and higher-forcing weather, but the diagnostic should not be reduced to a single climate-scenario ranking. Local climate, selected annual file, and building-zone response still shape the observed tail profile.

\subsection{PMV is a comparator, not a strawman}

The PMV results also support a restrained interpretation. PMV was correlated with $p_{\mathrm{tail}}$, and the standard $|\mathrm{PMV}|=0.5$ band was conservative in this fixed-reference panel. It almost never missed high-tail states. Its main difference from $p_{\mathrm{tail}}$ was over-flagging: many states outside the PMV neutral band did not cross the selected TSV tail-probability cutoff. The comparison is also geometrically informative. The learned expected-TSV scalar stays within a narrower empirical range and tracks the calibrated tail probability closely, while PMV spans a wider continuous comfort-index range and includes a lower-tail-probability branch at elevated $|\mathrm{PMV}|$. This tighter mapping is unsurprising because $\mu_{\mathrm{TSV}}$ and $p_{\mathrm{tail}}$ are two summaries of the same calibrated TSV probability vector; it should not be interpreted as a claim that the relationship is linear. That branch means that some states can look substantially displaced on the PMV scale while the calibrated TSV distribution still assigns most probability to neutral or mild bins. This is expected: $\mu_{\mathrm{TSV}}$ is a probability-weighted mean of the calibrated seven-class TSV predictor used in this study, whereas PMV is Fanger's continuous index built from thermal-balance terms and empirical chamber-vote evidence. The two scalars therefore stretch the same simulated thermal states differently.

Table~\ref{tab:pmv_scalar_comparison} makes the threshold consequence explicit. An offline optimized $|\mathrm{PMV}|$ threshold can be made closer to the $p_{\mathrm{tail}}$ screen, but it still disagrees more often than the optimized expected-TSV threshold. The conventional $|\mathrm{PMV}|=0.5$ band has a different purpose: it is protective as a comfort band, not tuned to reproduce this paper's tail-probability screen. The high false-positive rate under that band is therefore not a defect by itself; it shows that a conventional scalar comfort band and a calibrated severe-tail probability flag are not interchangeable diagnostics.

This means the diagnostic should not be framed as PMV failure. PMV remains a valuable standard comfort index and benchmark. The primary TSV model also uses PMV as one covariate, so the PMV panel is not an independent model contest. The no-PMV robustness check is therefore important, but it has a narrow meaning: the same stored PMV scalar was compared against two different TSV probability outputs, one from the primary model and one from a separately trained model without PMV. Removing PMV modestly loosened the PMV--$p_{\mathrm{tail}}$ relationship, which is consistent with PMV carrying useful information for reported TSV prediction, but it left the main aggregation result nearly unchanged. The probability diagnostic answers a different question: how much of the calibrated empirical TSV distribution lies in discomfort tails? For building studies that need a compact scalar, PMV or expected TSV may be enough. For studies that need tail-risk screening, the probability vector carries information that scalar summaries do not directly expose.

\subsection{Zone aggregation is a major source of information loss}

The zone results show that aggregation can matter more than the choice between two scalar summaries. Mean-zone aggregation gave 13.4\% modeled high-tail exceedance, while any-zone aggregation gave 37.0\%. This is not a minor reporting difference; it changes the apparent risk profile of the same fixed simulated building. The threshold sweep strengthens this point. Changing the high-tail screen changes the absolute exceedance percentage, as expected, but it does not remove the gap between mean-zone and any-zone summaries.

The zone pattern should not be summarized as a simple ``top floor is worst'' or ``bottom floor is worst'' result. In this prototype, P1 is the south perimeter zone and is the dominant local exceedance zone; bottom P1 was higher than middle and top P1 in most cases. At the same time, area-weighted floor exceedance was usually highest on the middle floor because the large core zone and the other perimeter orientations contributed more to the floor-level average. The any-zone floor metric answers a different question: whether at least one zone on that floor crosses the high-tail screen. It therefore exposes small-zone perimeter risk rather than suppressing it. This is why bottom-floor any-zone exceedance can be high even when the middle floor has the higher area-weighted mean. The conditioned top zones are also not direct roof-surface rooms; they sit below a top-floor plenum and roof assembly. The useful interpretation is therefore more specific: orientation, floor construction, zone size, aggregation rule, and local climate interact, and the probability diagnostic makes those interactions visible.

Figure~\ref{fig:zone_floor_position} also shows that floor differences are differences in degree, not separate thermal regimes. The bottom, middle, and top floor summaries are broadly comparable in magnitude, with the middle floor slightly elevated under area weighting. The tail-component bars show the same directional structure on all three floors: warm-tail probability dominates the discomfort signal, while cold-tail probability remains small. This supports the paper's diagnostic framing. The main finding is not that one floor is categorically unsafe, but that the conclusion changes depending on whether the analyst reports a mean, area-weighted, percentile, or any-zone view.

The result should still be read carefully. EnergyPlus zone states are well-mixed summaries, so the diagnostic does not prove that occupants at specific desks experienced radiant asymmetry, stratification, or local drafts. Any-zone exceedance is therefore a modeled zone-state screen, not a measured occupant-exposure rate. It shows that even at the modeled zone level, averaging can hide perimeter/core differences that are relevant to a probability-state diagnostic. A future implementation should report mean, percentile, and max-zone summaries together, or keep the probability object at the zone level when the building system can act by zone.

\subsection{Metabolic spread changes the risk profile}

The metabolic-spread diagnostic turns a common modeling convention into a measurable uncertainty. Holding the environment fixed while changing the metabolic input produced threshold flips in 19.0\% of mean-zone records and 38.7\% of max-zone records. This does not mean that EnergyPlus loads would be unchanged in a real population; the diagnostic deliberately did not rerun internal heat gains. It isolates the representation question: how sensitive is the predicted TSV tail probability to the representative metabolic assumption?

The answer is that the sensitivity is large enough to affect the risk profile. This is especially relevant because metabolic rate is often treated as a fixed office default. In the warm-dominant fixed-reference panel, lower metabolic-load assumptions generally lowered predicted tail exceedance because the same environmental state is interpreted with less occupant heat production. That is the model-output result. The demographic interpretation is more cautious. The metabolic-rate correction used to define the profiles shows that lower W/person values may correspond to population groups such as older or smaller-bodied occupants \citep{Guo2025Correcting120W}. The result should not be read as evidence that those occupants are safer or experience less actual discomfort. A recent review of elderly thermal characteristics emphasizes that older adults' reported thermal sensation and discomfort responses can differ from those of younger adults, including altered sensitivity and tolerance patterns \citep{Xu2025ElderlyThermalReview}. The present diagnostic therefore estimates the probability of reporting severe TSV categories under a trained TSV model; it does not measure latent discomfort, physiological heat strain, or health vulnerability. The non-monotone response between some adjacent profiles also cautions against overinterpreting the model as a smooth physiological curve. Plausible metabolic assumptions expose a spread in predicted reported-TSV tail risk that a single mean-occupant convention hides.

\subsection{Downstream use without overclaiming}

The probability object could support future operational designs. A fuzzy rule base could use calibrated tail probabilities as inputs. A stochastic or chance-constrained formulation could use $p_{\mathrm{tail}}$ or cumulative TSV probabilities as constraints. A building operator could use the same diagnostic to screen weather files, zones, or occupant-assumption scenarios before choosing an operating policy.

Those are downstream uses, not results of this paper. The fixed-reference simulations here do not test energy savings, comfort improvement, setpoint stability, equipment wear, demand response, or field deployment. The 15-minute records are diagnostic probability records. They should not be interpreted as equipment commands.

\section{Limitations and Future Work}\label{sec:limitations}

\paragraph{Static comfort calibration}
The TSV probability model is calibrated from contemporary ASHRAE Global Thermal Comfort Database II observations and is held fixed across all future-weather files. The study does not model future changes in expectations, clothing behavior, adaptive opportunity, physiology, culture, or thermal history. The results are therefore conditional stress-test diagnostics, not forecasts of how occupants will perceive buildings in the 2050s or 2080s.

\paragraph{Fixed building testbed}
The DOE Medium Office prototype is held fixed across all cases. This isolates the weather and representation effects, but it does not represent future envelope upgrades, code changes, adaptive glazing, shading retrofits, HVAC redesign, or future operation practice.

\paragraph{Weather-panel scope}
The 144 weather files are selected annual cases from a local CMIP-derived panel. They support comparisons across city, scenario, time slice, and severity, but they are not a probabilistic climate ensemble. The paper does not estimate event frequencies or regional climate uncertainty.

\paragraph{EnergyPlus zone resolution}
EnergyPlus provides well-mixed zone-level states. The zone-aggregation diagnostic therefore quantifies information loss between modeled zones. It does not resolve sub-zone radiant asymmetry, vertical stratification, local air speed, solar patches, or occupant-location-specific exposure. Reported any-zone exceedance rates should therefore be read as modeled zone-state exceedances rather than as measured occupant-exposure rates.

\paragraph{Tail calibration}
The discomfort tails are data-poor relative to the neutral class, especially at $\mathrm{TSV}=\pm3$. Post-hoc isotonic calibration improves probability reliability, but tail estimates remain more uncertain than central categories. Future work should test tail-weighted training, partial proportional-odds models, mixture-of-threshold models, and calibration methods that better preserve class dependence in sparse tails.

\paragraph{Metabolic-spread interpretation}
The metabolic diagnostic recomputes TSV probabilities over fixed environmental traces. It does not rerun EnergyPlus with altered occupant sensible heat gains, nor does it model correlations among metabolic rate, clothing, activity, age, sex, body size, or adaptation. The target is also reported TSV, not latent physiological strain or actual health risk. If some groups systematically report thermal sensation or discomfort differently from younger reference populations, or if TSV datasets encode reporting and adaptation biases, the diagnostic would inherit those biases. Lower predicted reported-TSV tail probability for a lower metabolic-rate profile should therefore not be equated with lower vulnerability. A full population model should represent these variables jointly rather than as independent scenario shocks and should distinguish reported sensation, comfort, preference, and physiological heat stress.

\paragraph{No closed-loop validation}
The present study intentionally excludes closed-loop operating validation. The diagnostic does not establish better HVAC performance, lower EUI, improved comfort, setpoint feasibility, equipment stability, or reduced demand. Those claims would require separate closed-loop experiments with explicit objectives, equipment constraints, dwell-time logic, and field or high-fidelity validation.

\section{Conclusions}

This paper presents a diagnostic-interface study of probabilistic thermal comfort risk. It is not a control study. The central object is the calibrated seven-class TSV distribution, from which expected TSV and discomfort-tail probability are both derived. The key definition is simple: $p_{\mathrm{tail}}=\Pr(\mathrm{TSV}\le-2)+\Pr(\mathrm{TSV}\ge+2)$, the predicted probability mass in the cold and warm discomfort tails.

In a 144-case fixed-reference Medium Office panel, 13.4\% of occupied records crossed $p_{\mathrm{tail}}\ge0.20$. High-tail exceedance increased in later time slices and under SSP5-8.5, showing how the fixed archetype appears under selected warming-climate weather stress. PMV and expected TSV both tracked tail risk, but neither replaces the probability distribution when the diagnostic question is explicitly about tail probability. Zone aggregation created the largest information loss: any-zone high-tail exceedance was 37.0\%, compared with 13.4\% under mean-zone aggregation. Metabolic-rate assumptions also changed threshold status, especially under max-zone aggregation.

The defensible conclusion is narrow and useful. Calibrated TSV probabilities reveal discomfort-tail, spatial, and occupant-assumption risk profiles that scalar comfort states compress. The result is a diagnostic representation that can inform future operational design, but it is not itself a validated operating strategy.
